\section*{Anexo 4-N — Correspondencia de componentes lógicos}
\phantomsection\label{fuente-15-anexo-4-1-n-correspondencia-de-componentes-logicos}
\phantomsection
\label{anx:N}
\addcontentsline{toc}{section}{Anexo 4-N — Correspondencia de componentes lógicos}
La matriz identifica los doce módulos y los servicios transversales para cotejarlos con el esquema y explicación de solución exigidos en 3.3--3.4 y con su realización física en 4.2. Las capacidades de 3.3 y 3.4 se corresponden con los módulos y contratos de 4.1. La matriz de correspondencia citada en 4.2.2 identifica su realización física.

La matriz completa vincula cada capacidad de 3.3–3.4 con su módulo o capa de 4.1, su recurso de 4.2 y su emplazamiento.

\begin{tablalafrox}{Correspondencia entre la arquitectura lógica y la física}{tab:mapeo-logica}%
  {F{0.23} F{0.20} Y F{0.20}}%
  {\cab{Capacidad de 3.3–3.4} & \cab{Módulo o capa (4.1)} & \cab{Componente físico (4.2)} & \cab{Dónde se ejecuta}}
 3.3.1–3.3.2, recepción en 3.4.1  &  M1 Recepción  &  A-01 y A-02, con C-03 y C-05  &  Sitio  \\
 3.3.1–3.3.3, FEFO, reserva y cross-docking en 3.4.1  &  M2 Inventario  &  A-01 y A-02 en el sitio, N-04 y N-05 para el consolidado  &  Sitio y nube  \\
 3.3.1 y 3.3.3, toma de pedido en 3.4.1  &  M3 Preventa  &  C-01 en el terminal, N-04, N-05 y N-07  &  Terreno y nube  \\
 3.3.1, planificación en 3.4.1  &  M4 Rutas  &  N-04 y N-05, motor de optimización en su propio contenedor  &  Nube  \\
 3.3.1–3.3.3, preparación y despacho en 3.4.1  &  M5 Preparación  &  A-01 y A-02, con C-03 y C-04  &  Sitio  \\
 3.3.1–3.3.3. Entrega en 3.4.1  &  M6 Reparto  &  C-02 en el terminal, N-04 y N-05, con la evidencia en S3  &  Terreno y nube  \\
 3.3.1, rendición y cierre en 3.4.1  &  M7 Cobranza y rendición  &  C-02, N-04, A-04 hacia el ERP  &  Terreno, nube y VM-04  \\
 3.3.1, devoluciones y cierre en 3.4.1  &  M8 Devoluciones y envases  &  C-02, N-04, A-01 para la recepción física de retornos mediante la puerta de API local  &  Terreno, nube y sitio  \\
  3.3.2, excursión y retiro sanitario en 3.4.1 & M9 Calidad y trazabilidad & B-01/B-02 en cámaras; gateway Greengrass de camión y TC58e con sensor independiente y termógrafo B-03; CloudConnect y N-08; N-04/N-06 para decisión y consulta & Sitio, camiones y nube \\
  3.3.1 y 3.3.3, indicadores y costo de servir en 3.4.1 & M10 Analítica & N-10 & Nube \\
 3.3.1 y 3.4.2, pedidos EDI  &  M11 Canal moderno  &  N-04 con el perfil EDI y el transporte AS2, A-04  &  Nube y VM-04  \\
 3.4.2, módulos y sistemas externos  &  M12 Telemetría  &  N-04 y N-06, con la API del proveedor de telemetría  &  Nube  \\
 3.4.2, quince actores e interfaces  &  Capa 1, portales y consolas  &  N-01 a N-03 y frontend de consolas en N-04  &  Nube  \\
 3.3.3 y 3.4.2, continuidad y accesos  &  Capas 2 y 3, entrada y puerta de enlace  &  CloudFront, WAF y API Gateway (N-01 a N-03 y N-12), puerta de API local en A-01  &  Nube y sitio  \\
 3.3.1–3.3.3, ERP, continuidad y datos conciliados  &  Capa 5, integración y eventos  &  A-03, N-09, A-04  &  Sitio y nube  \\
 3.3.1–3.3.3, lote, reserva, evidencia y consolidación  &  Capa 6, datos  &  A-02, N-05, N-06, N-07, N-10 y N-11  &  Sitio y nube  \\
 3.3.3 y 3.4.2, acceso por actor y turno  &  Capa 7, seguridad e identidad  &  A-05 con el verificador local, F-03, N-12  &  Nube y sitio  \\
 3.3.3, continuidad, incidentes y reconciliación  &  Capa 8, observabilidad  &  F-01 y CloudWatch (N-12)  &  Sitio y nube  \\
 3.4.3, implementación y promoción por ambientes & Cadena de entrega de 4.1.9 y 4.2.4.1 & DEV-PIPE y DEV-IAC, pipeline del CLIENTE, CodeBuild y ECR, perfiles firmados del artefacto & Cinco ambientes, con Desarrollo, QA y Preproducción aislados de los sitios \\
 3.4.4, implantación, convivencia y reversión & Olas de 4.1.8 y despliegue de 4.2.4.1 & A-01/A-02 por sitio, C-01/C-02/C-03 y versiones preservadas para reversión & Sitio, terreno y nube, escritor único por dominio \\
 3.4.5, operación, soporte y continuidad & Capa 8, identidad, gobierno de 4.2.4 a 4.2.6 & F-01/F-02/F-03 y N-12/N-13, gestión de incidentes y capacidad de 4-W.7 & Monitoreo central, colectores locales y soporte asistido \\
 3.4.6, participación y adopción por perfil & Capa 1, matriz de actores y permisos, olas de 4.1.8 & C-01/C-02/C-03, portales y consolas, QA de práctica & Terreno, sitios y nube, participación no implica permiso de acceso \\
\end{tablalafrox}

{\fontsize{9pt}{11pt}\selectfont\raggedright Fuente: (Bases Técnicas del caso, secciones~14.1--14.2 y~16.1, pp.~24--25 y~28--29), para volumetría y decisiones abiertas, y (Subdocumento 4, anexos, diseño de este apartado, pp.~\pageref*{fuente-15-anexo-4-1-n-correspondencia-de-componentes-logicos}--\pageref*{fuente-15-anexo-4-1-n-correspondencia-de-componentes-logicos-fin}).\par}

\begin{tablalafrox}{Correspondencia de módulos y realización física}{tab:correspondencia-modulos}{F{0.10} F{0.27} F{0.19} Y}{\cab{Módulo} & \cab{Capacidad en 3.3/3.4} & \cab{Contrato} & \cab{Realización en 4.2.2}}
M1 & Recepción & INT-03/06 & Perfil WMS local, base por sitio y ACL ERP. \\
M2 & Inventario & INT-03/04/12 & WMS por sitio, consolidación central y réplica de lectura de Talca. \\
M3 & Preventa & INT-01/06 & App Kotlin y API Laravel central. \\
M4 & Ruteo & INT-10 & Servicio de planificación y optimizador separado por contrato. \\
M5 & Picking y carga & INT-03/07 & HHT, puerta local y WMS con control de guía. \\
M6 & Entrega y POD & INT-02 & App de reparto, API y almacenamiento de evidencia. \\
M7 & Rendición y cartera & INT-02/06/09 & API, trabajadores Laravel y ACL de cobranza. \\
M8 & Devoluciones & INT-02/03/06 & Registro móvil/local y servicios de conciliación. \\
 M9  &  Trazabilidad y frío  &  INT-03/05 y API local autorizada  &  A-01/A-02/C-03: evaluación y disposición local. B-01/B-02 en CD y B-03 a B-15 con C-02 en camiones: captura y bloqueo, N-04/N-06/N-08: ingesta y consulta centrales.  \\
M10 & Analítica & Eventos canónicos & Ingesta a lago y servicios de consulta analítica. \\
M11 & Canal moderno & INT-08 & Reglas Laravel, OpenAS2 y conector por cadena. \\
M12 & Telemetría & INT-15 & Adaptador de solo lectura y consulta de ruta real. \\
\end{tablalafrox}
{\fontsize{9pt}{11pt}\selectfont\raggedright Fuente: (Bases Técnicas del caso, secciones~14.1--14.2 y~16.1, pp.~24--25 y~28--29), para volumetría y decisiones abiertas, y (Subdocumento 4, anexos, diseño de este apartado, pp.~\pageref*{fuente-15-anexo-4-1-n-correspondencia-de-componentes-logicos}--\pageref*{fuente-15-anexo-4-1-n-correspondencia-de-componentes-logicos-fin}).\par}

La tabla verifica la correspondencia de los doce módulos con sus capacidades, contratos y realizaciones. La Tabla~\ref{tab:mapeo-logica} citada en 4.2.2 completa el enlace físico. Las capacidades transversales también se corresponden: presentación y borde con aplicaciones, CloudFront y acceso privado. Puerta de enlace con API central y puerta local, eventos con RabbitMQ, shipper y SQS. Datos con bases, réplica de lectura y objetos, identidad con Keycloak y verificador local, secretos y auditoría con sus servicios de custodia. Observabilidad con OpenTelemetry/ADOT, alarma local y CloudWatch. Notificaciones se traza a INT-11. El motor de optimización de M4 se integra mediante un contrato separado y se ejecuta en su propio contenedor.

La correspondencia entre 3.3, 3.4, 4.1 y 4.2.2 se verifica en ambos sentidos mediante identificador, contrato, responsable y realización física.

\subsection*{Incorporación de las innovaciones del Subdocumento 13}
Las cinco innovaciones del Subdocumento 13 se realizan sobre los componentes existentes de la manera siguiente:
\begin{itemize}
\item Seguimiento de vencimiento posentrega en el local: C-01 y M3 reciben la precarga nocturna de M10/N-10 y registran la respuesta de la visita. M2 y M9 aportan los datos de inventario y trazabilidad. Las capas 1, 4, 5 y 6 soportan el recorrido. El saldo estimado en el local no modifica el inventario oficial ni autoriza disposición sanitaria.
\item Reproducción de incidentes de terreno: C-01/C-02/C-03 capturan el recorrido causal de M3/M6/M5 con límites de tamaño y frecuencia. OpenTelemetry y CloudWatch aportan la observación. Las capas 7 y 8 controlan acceso y trazas. Los recursos temporales se crean dentro de QA y Preproducción, no en un sexto ambiente, y se destruyen al finalizar, también ante error. Se excluyen nombres, credenciales, firmas, fotografías y contenido comercial. Al cerrar se elimina la captura y solo se conserva el escenario sintético de regresión. Operación recoge el caso, Calidad lo admite y Desarrollo corrige.
\item Vida útil remanente por historia térmica del lote: B-01/B-02 en cámaras y el kit B-03 a B-15 con C-02 en camiones alimentan N-06/N-08/N-10. B-04 evalúa las tres PT100 de B-05/B-06, activa B-12 y conserva la evidencia; C-02 evalúa por separado B-08 y contrasta el latido firmado del gateway. La protección local no depende de Internet. Móvil e Iridium SBD operan en paralelo con alarmas y órdenes de Calidad; productor, arranque y secuencia identifican un mismo evento recibido por ambos caminos. B-03 aporta el registro de aire completo al retorno, no es el único evaluador ni el único camino de aviso. M9 persiste la decisión y la obligación de aviso antes de A2. Las capas 1, 2, 4, 5, 6, 7 y 8 mantienen captura, protección, eventos, datos, custodia y observabilidad. B-02 acumula la historia de cámara durante 24 horas sin WAN. N-10 calcula el modelo determinista con parámetros aprobados por Calidad y conserva sus versiones. La innovación reutiliza los 30.548 mensajes y las 127.848 lecturas diarias de INT-05; no suma otra ingesta. No se rellenan tramos sin dato, no se extiende la fecha impresa, no se sustituye FEFO y la estimación no libera automáticamente un lote.
\item Tramo variable de la Operación ligado al costo de servir: M6/M7/M8 aportan hechos operacionales y M10/N-10 genera el informe mensual mediante las capas 4, 6 y 7. Finanzas del CLIENTE aprueba la liquidación. El ERP conserva la autoridad contable y recibe el resultado por la integración existente. El informe identifica período, índices, costo comparable, OTIF, recomendación y respuesta. Los montos pertenecen a la Oferta Económica.
\item Hoja de negocio del almacenero: M10/N-10 genera la hoja y M3/C-01 recibe la precarga cifrada. M6/C-02 registra su entrega e imprime en cabina. Las capas 1, 4, 6 y 7 limitan el acceso al receptor y distinguen copia pedida, entregada y en espera. El cliente no necesita aplicación, registro ni enlace para recibirla. El bloque agregado solo se habilita bajo las condiciones de privacidad del Subdocumento 13, con umbral mínimo de locales y supresión de celdas.
\end{itemize}

Las cinco innovaciones emplean los seis contratos de datos versionados del Subdocumento 13: aviso de vencimiento, respuesta de visita, paquete de reproducción, historia térmica, liquidación del tramo y hoja de negocio. Son intercambios sobre las superficies existentes, no seis contrapartes externas nuevas ni módulos adicionales. Sus productores y consumidores conservan autorización, idempotencia, retención por dominio y acuse durable cuando corresponda. La trazabilidad añade estos recorridos sin alterar el inventario de doce módulos, las 37 capacidades transversales ni las quince integraciones del dimensionamiento base.
\subsection*{Inventario trazable de capacidades transversales}
Los siguientes 37 identificadores complementan los doce módulos: el inventario lógico tiene 49 unidades verificables. La última columna identifica los recursos de 4.2 y T-11 que realizan cada capacidad, además de su límite de responsabilidad. La Tabla~\ref{tab:mapeo-logica} resume el mapeo por módulo y capa. Esta tabla permite verificar los componentes transversales por separado.
\begin{tablalafrox}{Inventario transversal y realización física}{tab:inventario-transversal}{F{0.16} F{0.17} F{0.17} F{0.15} Y}{\cab{ID} & \cab{Componente} & \cab{Capacidad} & \cab{Contrato} & \cab{Realización física y límite}}
 UI-MOV  &  Apps Kotlin / Room  &  preventa y entrega  &  INT-01/02  &  C-01/C-02 en terreno, N-04 en nube. M3/M6/M7/M8  \\
 UI-HHT  &  Terminales de bodega  &  preparación y disposición sanitaria autorizada  &  INT-03 / API local  &  C-03 y A-01 en los cinco sitios, perfiles separados para M1/M2/M5 y Calidad en M9.  \\
 UI-WEB  &  Portales Angular  &  Administración, portales  &  API autorizada  &  N-01 a N-03 y frontend N-04 en nube. Roles por empresa  \\
 BOR-PUB  &  CloudFront / AWS WAF / AWS Shield Advanced  &  seguridad  &  HTTPS  &  N-01 a N-03 y N-12 en nube. Entrada pública única  \\
 BOR-B2B  &  Superficie AS2  &  canal moderno  &  INT-08  &  N-04, OpenAS2 tras CloudFront, WAF y ALB interno en nube. Contrapartes registradas  \\
 GW-CENT  &  REST API / authorizer  &  plataforma  &  API / sincronización  &  N-04 y N-12, API Gateway y authorizer en nube. ADR-13  \\
 GW-LOCAL  &  Puerta API de sitio  &  WMS  &  INT-03  &  A-01 en VM-01, VM-C01 y E-01. Validación sin WAN  \\
 AUTH-CENT  &  Keycloak  &  seguridad  &  INT-13 / OIDC  &  A-05 en Fargate y N-05 en nube. Autoridad de identidad  \\
 AUTH-LOCAL  &  Verificador de turnos  &  continuidad  &  Manifiesto firmado  &  A-05 en VM-05, VM-C03 y E-01. PIN, turno y sitio  \\
 WMS-SITE  &  Perfil WMS local  &  WMS  &  INT-03  &  A-01/A-02 en VM-01/02, VM-C01/02 y E-01. Un escritor por sitio  \\
 INT-BROKER  &  RabbitMQ  &  integración  &  AMQP  &  A-03 en VM-03, VM-C04 y E-01. Persistencia y DLQ  \\
 INT-SHIPPER  &  Shipper local a nube  &  integración  &  JSON / SQS  &  A-01 en VM-03, VM-C04 y E-01. N-09. Conserva outbox hasta acuse de recepción, confirmación de negocio exige resultado durable del consumidor.  \\
 INT-CONSUMER  &  Consumidor PHP  &  integración  &  AsyncAPI / JSON  &  N-04 con acceso a N-05/N-09. Persistir antes de acuse  \\
 INT-JOBS  &  Trabajadores Laravel  &  notificaciones y EDI  &  Colas internas  &  N-04/N-09 en nube. Perfil de trabajos en N-04 (Fargate)  \\
 INT-SCHED  &  Planificador  &  administración  &  Tareas versionadas  &  N-04 en nube. Un líder por ambiente  \\
 INT-ACL  &  Adaptador ERP  &  ERP / DTE  &  INT-06/07  &  A-04 en VM-04 de Talca. ERP único emisor  \\
 INT-EDI  &  OpenAS2 / perfiles  &  canal moderno  &  INT-08  &  N-04 en nube. A-04 para intercambio ERP. MDN no es acuse comercial  \\
 INT-GIS  &  Adaptador de mapas  &  rutas  &  INT-10  &  N-04 y contraparte GIS de INT-10. Tiempo máximo / precarga  \\
 INT-OPT  &  Optimizador M4  &  rutas  &  Trabajo versionado  &  Contenedor propio N-04 en nube. Restricciones de ruta  \\
 INT-GPS  &  Adaptador de flota  &  telemetría  &  INT-15  &  N-04/N-06 y contraparte de INT-15. Solo lectura  \\
 INT-NOTIF  &  Notificaciones  &  administración  &  INT-11  &  N-04 y SES/SNS/End User Messaging Social de INT-11. Vigencia por canal  \\
 INT-IOT  &  Captura térmica local  &  frío  &  INT-05  &  B-01/B-02 en CD; B-03 a B-15 y C-02 en camiones. B-04 y C-02 evalúan localmente; móvil/Iridium hacia N-08, sin duplicar el evento  \\
 DAT-OLTP  &  PostgreSQL / PostGIS  &  inventario  &  Repositorios propios  &  A-02 por sitio, N-05 en nube. Autoridad por sitio  \\
 DAT-CDC  &  Réplica de Talca  &  continuidad  &  INT-12  &  VM-02 y DMS/Aurora de N-05. Lectura, no doble escritor  \\
 DAT-OLAP  &  DynamoDB / Glue / S3 / Redshift  &  analítica  &  ETL / eventos  &  N-06/N-10 en nube, posiciones excluidas de us-east-1. Separado de OLTP  \\
 DAT-OBJ  &  Objetos POD / DTE  &  entrega  &  Hash y objeto  &  S3 de N-11 en nube. Retención por dominio  \\
 DAT-CACHE  &  Valkey central  &  plataforma  &  Lecturas autorizadas  &  N-07 en nube. No reserva stock  \\
 DAT-LOCAL  &  Room / cola cifrada  &  movilidad  &  UUID / sincronización  &  SQLite/Room en C-01/C-02/C-03. Separar fotos y registros sin confirmar  \\
 SEC-SECRETS  &  Custodia de claves  &  seguridad  &  Identidad de servicio  &  KMS, Secrets Manager y Parameter Store de N-12. Rotación segregada  \\
 SEC-AUDIT  &  Auditoría inalterable  &  administración  &  Evento de decisión  &  S3 Object Lock de N-11 y CloudTrail de N-12, bitácora local A-02. Sujeto, regla y resultado  \\
 SEC-SIEM  &  Detección y correlación  &  seguridad  &  Alertas  &  F-03 y servicios de detección de N-12. Casos Q/R  \\
 OBS-COLLECT  &  OpenTelemetry / ADOT  &  observabilidad  &  INT-14 / OTLP  &  F-01 en VM-06, VM-C04 y E-01. Buffer de 24 horas  \\
 OBS-LOCAL  &  Alarmas del sitio y del camión  &  continuidad  &  Alerta local  &  B-02 en CD, B-04/B-12 y C-02/B-08 en camión, y colectores F-01. No depende de WAN  \\
 OBS-CENT  &  CloudWatch  &  observabilidad  &  INT-14  &  CloudWatch de N-12 en nube. Plataforma única  \\
 DEV-PIPE  &  CI/CD y SBOM  &  plataforma  &  Artefacto firmado  &  Pipeline del CLIENTE, CodeBuild/ECR de N-04. T-11, entrega continua. Promoción por contrato  \\
 DEV-IAC  &  Terraform / Ansible  &  infraestructura  &  Estado y configuración  &  F-02, repositorio del CLIENTE y cuentas de los cinco ambientes. Propietario único  \\
 DEV-MDM  &  Gestión de terminales  &  seguridad  &  Enrolamiento  &  N-13 en nube y agentes en C-01/C-02/C-03. Política y borrado  \\
\end{tablalafrox}

{\fontsize{9pt}{11pt}\selectfont\raggedright Fuente: (Bases Técnicas del caso, secciones~14.1--14.2 y~16.1, pp.~24--25 y~28--29), para volumetría y decisiones abiertas, y (Subdocumento 4, anexos, diseño de este apartado, pp.~\pageref*{fuente-15-anexo-4-1-n-correspondencia-de-componentes-logicos}--\pageref*{fuente-15-anexo-4-1-n-correspondencia-de-componentes-logicos-fin}).\par}
\subsection*{Correspondencia de actores y permisos del sistema}
La Tabla~\ref{tab:actores-permisos} realiza el catálogo de quince actores del apartado 3.4.2 del Subdocumento 3. Cada acción se autoriza por separado sobre el ámbito indicado. Ni el cargo ni el acceso a una consola conceden todos los permisos de un módulo.

\begin{tablalafrox}{Actores del sistema, interfaces y permisos}{tab:actores-permisos}{F{0.17} F{0.22} F{0.24} F{0.24} Y}{\cab{Actor} & \cab{Interfaz y módulos} & \cab{Acciones} & \cab{Ámbito y límite} & \cab{Etapa}}
Preventista & App de preventa, M3, lectura M2/M7 & Consultar stock y crédito, capturar y sincronizar pedidos & Su cartera, no reserva offline & 1 \\
Conductor propio & App de reparto, M6/M7/M8 & Entrega, POD, cobro, retorno y envases & Su ruta y turno, no aprueba su descuadre & 1 \\
Conductor externo & App de reparto, M6/M7/M8 & Entrega y rendición propias & Identidad personal ligada a empresa y turno & 1 \\
Preparador & Terminal de bodega, M5, M1/M2 con permiso & Lecturas, faltantes y movimientos autorizados & Sitio, misión y turno, no libera lotes & 1 \\
Cliente del canal tradicional & Portal opcional, M3/M6/M7 & Armar pedidos, ver entregas, documentos y saldo & Solo sus datos, compra asistida sin cuenta & 2 \\
Cliente del canal moderno & Portal y conector EDI, M11/M3 & Pedidos, consultas y mensajes EDI & Su cadena, sesión humana y técnica separadas & 2 \\
Empresa transportista & Portal por representante, M4/M6/M12 & Ver rutas y documentos, confirmar conductor y vehículo & Su empresa, no firma POD & 2, registro por despacho en 1 \\
Proveedor & Portal por representante, M1 & Consultar órdenes y recepciones & Su organización, solo lectura & 2 \\
 Jefa de Calidad  &  Consola Angular central y perfil Kotlin de Calidad por puerta local, M9, acciones M2/M5  &  En conexión: fijar umbral y duración. En el sitio: evaluar, liberar, bloquear o retirar con evidencia  &  Lote e instalación, decisión exclusiva y permiso personal vigente, sin cambiar maestros offline  &  1  \\
Gerente Comercial & Consola comercial, M3/M7/M10/M11 & Reglas, excepciones e indicadores & Política y canal, con historial & 1, canal moderno en 2 \\
Gerente de Finanzas & Consola financiera, M7/M10 & Supervisar rendición y reglas financieras & Tesorería nominada. Quien registra no aprueba & 1, costo en 2 \\
Planificador de Rutas & Consola de rutas, M4/M12 & Generar, corregir y aprobar rutas & Motivo auditado, no altera cobros & 1 \\
Jefe de TI & Consola de administración & Identidad, configuración, integraciones y monitoreo & Privilegio temporal con MFA, sin permisos de negocio & 1 \\
Gerente de Operaciones & Consola operacional, M2/M4/M5/M6 & Supervisar despacho y excepciones & Sitio, no omite guía ni Calidad & 1 \\
Jefa de Bodega & Consola de bodega, M1/M2/M5/M8 & Recepción, inventario, FEFO, preparación y retornos & Sitio, ajustes justificados & 1 \\
\end{tablalafrox}
{\fontsize{9pt}{11pt}\selectfont\raggedright Fuente: elaboración a partir del apartado 3.4.2 y del Anexo 3.I del Subdocumento 3. Diseño propuesto: (Subdocumento 4, anexos, diseño de este apartado, pp.~\pageref*{fuente-15-anexo-4-1-n-correspondencia-de-componentes-logicos}--\pageref*{fuente-15-anexo-4-1-n-correspondencia-de-componentes-logicos-fin}).\par}

Recepción, despacho, catálogo, abastecimiento y Tesorería son funciones de personas nominadas con permisos específicos en los perfiles del CLIENTE. Abastecimiento aprueba la sugerencia de reposición de M2 sin crear otro actor. AL-ACT-01 del Anexo 4-V prueba cada fila con una operación permitida y otra denegada.

\subsection*{Comportamiento transversal de interfaces}
\label{anxN:interaccion}
Este catálogo concentra los criterios de interacción de la capa de presentación. Cada función mantiene el módulo dueño, el permiso y la prueba de accesibilidad; AL-UI-01 del Anexo 4-V define la evidencia de aceptación.

Cada perfil de los portales, consolas Angular y aplicaciones Kotlin tiene accesos a todas sus funcionalidades principales desde su pantalla de inicio en un máximo de tres interacciones. Este límite se comprueba por función y perfil en los prototipos y en la aceptación (RT-13.05, Bases Técnicas Transversales, cap. 13, p. 26).

Angular/Tailwind y Kotlin comparten un sistema de diseño documentado: paleta de hasta cinco colores principales, jerarquía tipográfica de a lo más dos familias, iconografía coherente, retícula y componentes reutilizables con sus estados y reglas de uso (RT-13.09, Bases Técnicas Transversales, cap. 13, p. 26).

Los portales, consolas y aplicaciones permiten elegir modo claro u oscuro, tamaño de texto y densidad de los listados, conservando la preferencia de cada persona en su dispositivo y manteniendo la accesibilidad. Se entregan en español. El caso no exige otros idiomas para sus perfiles operacionales (RT-13.12, Bases Técnicas Transversales, cap. 13, p. 26).

El Portal de Clientes ofrece una vista ligera para teléfonos de bajo costo o de generaciones anteriores, con imágenes bajo demanda, listados paginados y sin animaciones, conservando pedido, estado de entrega, documentos y saldo. Usa los navegadores y sistemas compatibles de la \href{LAFROX-Subdocumento4-Anexos.pdf#tab:inventario-software}{Tabla de inventario complementario} del Anexo 4-P y se prueba en equipos de esas categorías antes de su aceptación (RT-17.08, Bases Técnicas Transversales, cap. 17, p. 32).

Los portales y la consola web en Angular 22 con Tailwind CSS adaptan su contenido a escritorio, tableta y teléfono conforme a RT-13.02 (Bases Técnicas Transversales, cap. 13, p. 26). El diseño parte de la pantalla pequeña y utiliza los puntos de quiebre de Tailwind: \texttt{sm} 40 rem, \texttt{md} 48 rem, \texttt{lg} 64 rem, \texttt{xl} 80 rem y \texttt{2xl} 96 rem, equivalentes a 640, 768, 1.024, 1.280 y 1.536 píxeles con la referencia de 16 píxeles por rem (Tailwind Labs, s.~f.). En teléfono, la navegación se recoge en un menú, los formularios se ordenan en una columna y los listados presentan tarjetas con acceso al detalle. En tableta se distribuyen los controles y en escritorio se muestran paneles y columnas adicionales. Se conservan todas las operaciones y los datos necesarios, sin reducir simplemente la escala. La aceptación comprueba cada transición y ambas orientaciones. Los terminales Zebra con Kotlin nativo usan la pantalla fija de cada modelo y un diseño específico para su perfil, con tamaños legibles y controles accesibles.

La matriz de navegadores, versiones y actualización de la \href{LAFROX-Subdocumento4-Anexos.pdf#tab:inventario-software}{Tabla de inventario complementario} del Anexo 4-P se aplica a los portales y a la consola web (RT-13.10, Bases Técnicas Transversales, cap. 13, p. 26). Las versiones exactas se registran antes de cada liberación y se contrastan con el conjunto Baseline de la versión de Angular utilizada (Angular, 2026b). El soporte del framework no sustituye las pruebas de la aplicación.

Todas las interfaces permiten recorrer y ejecutar íntegramente sus funciones por teclado, con orden de foco acorde con la lectura y la secuencia de trabajo, foco visible, acceso directo al contenido y ausencia de bloqueos del foco. Los diálogos reciben el foco al abrirse y lo devuelven al control de origen al cerrarse. Tablas, menús y formularios tienen controles semánticos operables mediante Tab, Shift+Tab, flechas, Enter y Escape según corresponda. Se ofrecen atajos documentados y desactivables para buscar, guardar un borrador y volver al inicio del perfil, sin interferir con los del navegador, del sistema ni con la escritura en campos. En Kotlin se verifica además la navegación con teclado físico o conectado al terminal. La aceptación recorre cada flujo por teclado, comprueba el orden y la recuperación del foco y conserva la conformidad WCAG 2.2 AA comprometida (World Wide Web Consortium [W3C], 2023; RT-13.11, Bases Técnicas Transversales, cap. 13, p. 26).

En el canal moderno, el cliente puede registrarse por sí mismo en el Portal de Clientes con el RUT de la empresa, un correo corporativo y una contraseña administrada por Keycloak. Antes de activar la cuenta, la integración con el ERP verifica que el RUT corresponda a un cliente activo habilitado por Puelche. El portal envía un correo de confirmación. Si el RUT no está habilitado, indica que se debe contactar al ejecutivo comercial. En el canal tradicional, la activación asistida por SMS sigue siendo opcional (RF-12.25, Subdocumento 3, Anexo 3.A).

LafroX verifica la conformidad con WCAG 2.2 nivel AA en todas las interfaces de uso humano, con herramientas automáticas y con pruebas manuales de los flujos, incluidos el teclado y las tecnologías de asistencia. Todo incumplimiento confirmado de nivel A o AA impide la promoción a Preproducción y la aceptación hasta que se corrija. El informe de conformidad y la evidencia de las correcciones se entregan al CLIENTE junto con los expedientes de aceptación (RT-13.01, Bases Técnicas Transversales, cap. 13, p. 26).

Cada acción muestra un resultado visual claro, que distingue entre operación confirmada, pendiente de sincronización y rechazada. Los mensajes de error explican en lenguaje simple qué ocurrió y qué debe hacer el usuario, sin mostrar códigos ni trazas técnicas (RT-13.06, Bases Técnicas Transversales, cap. 13, p. 26).

Las interfaces ofrecen alto contraste, controles táctiles de al menos 44 × 44 píxeles, iconos acompañados de texto y flujos guiados paso a paso para usuarios con baja alfabetización digital. En picking se aplican además el tamaño mínimo de 15 × 15 mm y la ausencia de gestos complejos descritos en 4.2.1.2 (RT-13.07, Bases Técnicas Transversales, cap. 13, p. 26).

La pantalla de inicio de cada usuario reúne en una bandeja única las aprobaciones, revisiones y excepciones pendientes de todos los módulos, según sus permisos. En la aplicación Kotlin incluye también los pendientes locales registrados durante un corte. Las tareas se ordenan por criticidad y plazo, con alertas de vencimiento, y cada una se resuelve en el flujo del módulo responsable (RT-16.13, Bases Técnicas Transversales, cap. 16, p. 30).

La búsqueda global de portales, consolas y perfiles Kotlin consulta las interfaces públicas de los módulos. Usa índices de texto completo en PostgreSQL, tolera errores de escritura y ofrece filtros por tipo de información. Cada módulo entrega solo los resultados que el usuario puede ver, también en los conteos de cada filtro. La aceptación verifica una respuesta de hasta dos segundos (RF-17.11). Durante un corte, la aplicación Kotlin busca solo en los datos locales autorizados (RT-16.27, Bases Técnicas Transversales, cap. 16, p. 31).

Los listados de portales, consolas y perfiles Kotlin se pueden ordenar, filtrar y paginar dentro de los permisos del usuario. La exportación en CSV UTF-8 o JSON respeta los filtros y el orden elegidos e incluye todos los registros filtrados, no solo la página visible. Si la exportación necesita datos centrales y no hay enlace, queda pendiente hasta que se recupere la conexión (RT-16.28, Bases Técnicas Transversales, cap. 16, p. 31).

\label{fuente-15-anexo-4-1-n-correspondencia-de-componentes-logicos-fin}
